
\begin{spacing}{1.2}
      \chapter{PENDAHULUAN}
\end{spacing}


\pagenumbering{arabic}
\vspace{4ex}

\section{Latar Belakang}
Penyu merupakan salah satu spesies yang dilindungi karena populasinya terancam punah di seluruh dunia.
Indonesia sendiri menjadi habitat penting bagi penyu, karena terdapat enam dari tujuh spesies penyu dunia yang hidup di perairan Indonesia \cite{JKT602}.
Namun, dalam beberapa dekade terakhir, populasi penyu mengalami penurunan yang signifikan.

Di alam liar, tukik (penyu yang baru menetas) menghadapi berbagai ancaman alami seperti predasi oleh burung, kepiting, dan biawak.
Meskipun demikian, ancaman terbesar terhadap kelangsungan hidup penyu berasal dari aktivitas manusia.
Pembangunan pesisir yang tidak terkendali menyebabkan berkurangnya habitat peneluran.
Selain itu, praktik perburuan ilegal untuk mengambil telur, daging, kulit, dan cangkang penyu turut mempercepat penurunan populasi secara global \cite{JKT602}.

Beberapa kasus kematian penyu akibat aktivitas ilegal telah dilaporkan di Sulawesi Barat.
Pada tahun 2016 ditemukan dua ekor penyu mati di Pantai Palippis dan Pantai Lapeo dengan indikasi pengambilan telur secara paksa.
Pada tahun 2017, tiga ekor penyu ditemukan dalam kondisi tubuh tidak utuh di Pantai Mampie.
Selanjutnya pada tahun 2019, sebanyak 18 ekor penyu ditemukan mati di lokasi yang sama dan diduga dibunuh untuk diambil telur serta dagingnya \cite{nur2022pelatihan}.
Kasus-kasus tersebut menunjukkan urgensi upaya konservasi yang lebih efektif.

Salah satu pendekatan konservasi yang penting adalah mempelajari perilaku dan pola migrasi penyu melalui pelacakan individu di habitat alaminya.
Teknologi yang umum digunakan saat ini adalah \textit{satellite tagging}, yaitu metode pelacakan posisi menggunakan modul GPS berpresisi tinggi yang dipasang pada tubuh penyu.
Teknologi ini memungkinkan peneliti memantau pergerakan, habitat, dan jalur migrasi penyu secara akurat.
Namun, metode ini memiliki beberapa keterbatasan, antara lain biaya perangkat yang tinggi, risiko kerusakan atau malfungsi perangkat, serta kebutuhan pemasangan yang tepat agar tidak membahayakan penyu \cite{10.3389/fmars.2018.00432}.

Oleh karena itu, diperlukan metode alternatif yang lebih efisien dan non-invasif untuk mengidentifikasi serta melacak individu penyu.
Salah satu pendekatan yang menjanjikan adalah re-identifikasi (Re-ID) berbasis visi komputer.
Re-ID merupakan proses mengenali kembali individu yang sama pada citra yang berbeda dengan memanfaatkan karakteristik visual unik \cite{zheng2016personreidentificationpastpresent}.
Pendekatan ini memungkinkan identifikasi individu meskipun terdapat variasi sudut pandang, pencahayaan, dan kondisi lingkungan.

Perkembangan teknologi \textit{deep learning}, khususnya \textit{Vision Transformer} (ViT), menunjukkan kemampuan yang sangat baik dalam mengekstraksi fitur visual yang kompleks.
Dibandingkan dengan arsitektur Convolutional Neural Network (CNN), ViT mampu menangkap hubungan global pada citra sehingga berpotensi menghasilkan performa identifikasi yang lebih tinggi \cite{9716741}.
Dengan memanfaatkan pola unik pada bagian kepala penyu, teknologi ini dapat digunakan untuk mengidentifikasi individu penyu secara otomatis dan non-invasif.

\section{Rumusan Masalah}
Berdasarkan latar belakang di atas, maka permasalahan yang dapat dirumuskan dalam penelitian ini adalah sebagai berikut:
\begin{enumerate}
      \item Bagaimana performa model Re-ID penyu berbasis ViT dengan pembelajaran metrik pada pola kepala penyu.
      \item Faktor apa saja yang mempengaruhi performa model dalam re-identifikasi citra penyu?
      \item Seberapa baik kemampuan model dalam generalisasi data yang memiliki pengambilan citra yang berbeda?
\end{enumerate}

\section{Tujuan}
\begin{enumerate}
      \item Menganalisis performa model Re-Identification (Re-ID) penyu berbasis Vision Transformer (ViT) dengan pendekatan metric learning dalam mengenali individu penyu berdasarkan pola pada bagian kepala.
      \item Mengidentifikasi dan menganalisis faktor-faktor yang memengaruhi performa model dalam proses re-identifikasi citra penyu.
      \item Mengevaluasi kemampuan generalisasi model dalam melakukan re-identifikasi pada data citra penyu yang diperoleh dengan variasi kondisi pengambilan gambar, seperti sudut pengambilan, pencahayaan, kualitas citra, dan perangkat yang digunakan.
\end{enumerate}

\section{Batasan Masalah}
Penelitian ini memiliki batasan masalah untuk menetapkan fokus penelitian. Batasan masalah meliputi:
\begin{enumerate}
      \item Penelitian berfokus pada pengembangan model ViT untuk Re-ID penyu, dengan arsitektur pra-latih DINOv3.
      \item Identifikasi hanya menggunakan citra kepala penyu.
      \item Citra masukan diasumsikan telah ter-crop pada region kepala; deteksi dan segmentasi kepala berada di luar lingkup penelitian.
      \item Penelitian tidak membahas aspek implementasi perangkat keras lapangan maupun integrasi sistem konservasi secara langsung.
\end{enumerate}

\section{Manfaat}
Penelitian ini diharapkan dapat memberikan beberapa manfaat sebagai berikut:
\begin{enumerate}
      \item Berkontribusi penelitian di bidang computer vision untuk konservasi satwa liar.
      \item Menjadi referensi bagi penelitian lanjutan terkait animal Re-ID berbasis Vision Transformer.
      \item Menyediakan opsi untuk identifikasi penyu laut yang non-invasif.
      \item Menyediakan peta sensitivitas parameter yang dapat dijadikan titik awal penyetelan bagi penelitian Re-ID satwa berbasis ViT berikutnya.
\end{enumerate}